\documentclass[10pt, dvipsnames]{beamer}

\usecolortheme{dove}
\usetheme{boxes}

\setbeamercolor{footline}{fg=black}
\setbeamerfont{footline}{series=\bfseries}

\usepackage[utf8]{inputenc}

\usepackage{hyperref}
\usepackage{enumerate}
\usepackage{subcaption}
\usepackage{graphicx}
\usepackage{animate}
\usepackage{hyperref}
\usepackage{mathtools}
\usepackage{amsmath}
\usepackage{indentfirst}
\usepackage{xspace}
\usepackage{tcolorbox}

\beamertemplatenavigationsymbolsempty

% ---------- Bibliography (biblatex+biber) ----------
\usepackage[
  backend=biber,
  style=phys,        % if missing on TeX Live 2022, use style=numeric
  sorting=ydnt,      % year (desc), name, title
  maxbibnames=3,     % short author lists for CV
  minbibnames=1,
  giveninits=true,   % initials instead of full given names
  defernumbers=true, % number only printed entries (no gaps)
  doi=true,
  url=false,         % hide URLs if DOIs present
  eprint=true
]{biblatex}

\date{\textcolor{gray}{2026}}

\setbeamertemplate{footline}{
{%
   \hspace{.95cm} \begin{minipage}{0.1\paperwidth} \insertdate \hfill \break \end{minipage} \hfill 
   \includegraphics[width=0.14\textwidth]{logos/phenix_small.pdf} 
   \hfill \begin{minipage}{0.045\paperwidth} \textcolor{gray}{\insertframenumber/\inserttotalframenumber} \hfill \break \end{minipage} \unskip \hspace{0.95cm}
   \vfill
}}

\setbeamertemplate{frametitle}[default][center]
\setbeamerfont{frametitle}{series=\bfseries}

\setbeamertemplate{caption}[numbered]

\AtBeginPart{\frame{\partpage}}

\newcommand{\etat}{\mbox{$\eta$}\xspace}
\newcommand{\cuau}{\text{Cu}+\text{Au}}
\newcommand{\pit}{\mbox{$\pi^0$}\xspace}
\newcommand{\Gevc}{GeV/$c$\xspace}
\newcommand{\Gevcsqr}{GeV/$c^2$\xspace}
\newcommand{\sqsn}{\sqrt{s_{_{NN}}}}
\newcommand{\pt}{\mbox{$p_T$}\xspace}
\newcommand{\norbertnote}{AN1270\xspace}
\newcommand{\zharkonote}{AN1281\xspace}

\newcommand{\photon}{\upgamma}


\begin{document}

\begin{frame}{PHENIX Preliminary Data Request \\ Low-$p_T$ Direct Photons}
   Name: Tongzhou Guo, A. Drees, M. Mitrankova\break
   PWG: Heavy Ion, Photon, PLHF\break
   What Observable: low-$p_T$ direct photon production using the external conversion method; $R_\gamma$, direct photon invariant yield, and integrated direct photon yield\break
   System, Energy, Run: Run-12 Cu$+$Au at $\sqrt{s_{NN}}=200$ GeV\break
   Last PWG Presentation Link:\break
   \textcolor{RoyalBlue}{\url{https://phenix-intra.sdcc.bnl.gov/cdsagenda/fullAgenda.php?ida=a2626}}\break
   Analysis Note: an1525 \break
   Analysis Note Web Link:\break
   \textcolor{RoyalBlue}{\url{https://phenix-intra.sdcc.bnl.gov/phenix/WWW/p/draft/tongzhouguo/Direct_photon_production_Cu_Au_v1.pdf}}\break
\end{frame}

\begin{frame}{Plots requested for preliminary status}
   \scriptsize
   \begin{itemize}
   \item[$\bullet$] $R_\gamma$ in Cu$+$Au minimum bias compared with PHENIX Au$+$Au 40--50\%.
   \item[$\bullet$] $R_\gamma$ in Cu$+$Au minimum bias compared with PHENIX Au$+$Au 50--60\%, 40--50\%, and 30--40\%.
   \item[$\bullet$] Direct photon invariant yield in Cu$+$Au minimum bias compared with PHENIX Au$+$Au 40--50\%.
   \item[$\bullet$] Integrated direct photon yield $dN_{\gamma}^{\rm dir}/dy$ for $1<p_T<5$ GeV/$c$ versus $dN_{\rm ch}/d\eta$.
   \end{itemize}

   \vspace{0.1em}

   \begin{center}
   \begin{tabular}{cc}
   \includegraphics[width=0.38\textwidth]{Pictures/rgamma_stat_syst0.png}
   &
   \includegraphics[width=0.38\textwidth]{Pictures/rgamma_stat_syst1.png}
   \\
   \includegraphics[width=0.38\textwidth]{Pictures/spectrum_stat_syst.png}
   &
   \includegraphics[width=0.38\textwidth]{Pictures/dndy_vs_dNch.png}
   \end{tabular}
   \end{center}
\end{frame}


\begin{frame}{Basic Information \& Provenance}
\textbf{Data set:} \\
Run-12 Cu$+$Au at $\sqrt{s_{NN}}=200$ GeV, minimum-bias data sample \\
\vspace{0.4em}

\textbf{Data provenance and reconstruction:} \\
CNT, DST\_EVE, pro99 \\
\vspace{0.4em}

\textbf{Analysis code location:} \\
\texttt{offline/AnalysisTrain/DileptonAnalysis/} \\
\texttt{offline/AnalysisTrain/pat/macro/Run\_DileptonAnalysis\_CuAu.C} \\
\vspace{0.4em}

\textbf{Please list all recalibrators used in this analysis:} \\
BeamOffset, MomScale, EMCal \\
\vspace{0.4em}

\textbf{Were any of these not in the MasterRecalibrator?} \\
No \\
\vspace{0.4em}

\textbf{Basic analysis cuts used:} next slide
\end{frame}


\begin{frame}{Basic Analysis Cuts}
   \small
   The nominal cut definitions are taken from the analysis note.  The main selections are:
   \begin{itemize}
      \item Event selection: BBCLL1($>1$ tubes)\_narrowvtx trigger and $|z_{\rm vtx}|<10$ cm.
      \item Charged tracks: quality $=31|51|63$, $|zed|<75$ cm, and $p_T>0.3$ GeV/$c$.
      \item Electron ID: RICH requirement $n_0\geq 2$, EMCal matching $|emcdz|<20$ cm and $|emcd\varphi|<0.05$ rad, and $E/p>0.6$ requirement.
      \item EMCal photon candidates: cluster $E_{\rm core}$ threshold, shower-shape requirement $\chi^2<3$, EMCal dead maps, and an electron/positron veto.
      \item Conversion pairs: external conversion-pair cuts based on variables such as $dzed<4$ cm, $1<r_{\rm conv}<29$ cm.
   \end{itemize}
   \vspace{0.4em}
   The systematic study repeats the extraction with controlled variations of these selections, including conversion-pair cuts, electron ID, electron-track momentum threshold, photon threshold, shower shape, energy scale, and $\pi^0$ peak extraction.
\end{frame}
\begin{frame}{What's the observable}
   \textbf{What observable is analyzed?}
   \begin{itemize}
      \item[$\bullet$] Low-$p_T$ direct photons in Cu$+$Au minbias at $\sqrt{s_{NN}}=200$ GeV
      \item[$\bullet$] Direct photon double ratio $R_\gamma$
      \item[$\bullet$] Direct photon invariant yield
      \item[$\bullet$] Integrated direct photon yield $dN_{\gamma}^{\rm dir}/dy$ for $1<p_T<5$ GeV/$c$
   \end{itemize}

   \textbf{What will be plotted?}
   \small
   \begin{itemize}
      \item[$\bullet$] $R_\gamma$ in 093 Cu$+$Au compared to PHENIX Au$+$Au 40--50\% from PPG243
      \item[$\bullet$] $R_\gamma$ in 093 Cu$+$Au compared to PHENIX Au$+$Au 50--60\%, 40--50\%, and 30--40\% from PPG243
      \item[$\bullet$] Direct photon invariant yield in Cu$+$Au minimum bias compared to PHENIX Au$+$Au 40--50\%
      \item[$\bullet$] Integrated direct photon yield $dN_{\gamma}^{\rm dir}/dy$ vs. $dN_{\rm ch}/d\eta$ compared to PHENIX and ALICE measurements
   \end{itemize}

   \textbf{What's new about this result?}
   \begin{itemize}
      \item[$\bullet$] First low-$p_T$ direct photon result in Run-12 Cu$+$Au collisions 
   \end{itemize}

   \textbf{What is the question that this new result will answer?}
   \begin{itemize}
      \item[$\bullet$] Whether Cu$+$Au direct photon production follows the same trend with $dN_{\rm ch}/d\eta$ as Au$+$Au, Cu$+$Cu, and Pb$+$Pb systems
   \end{itemize}
\end{frame}

\begin{frame}{Outline of the analysis method}
\begin{itemize}
\item[$\bullet$] If you use an established analysis method that was previously used in PHENIX published paper(s), please list the PHENIX paper(s) that uses the same method: PPG243

  \begin{minipage}{0.58\textwidth}
  \item[$\bullet$] If the method is used in previously approved preliminary results, please give the information on the approval, i.e., when approved and some of the approved plot:
  \begin{itemize}
      \item Au$+$Au direct photons using the external conversion method
      \item Preliminary granted to Wenqing Fan on 2019-10-28
      \item Supporting analysis note: an1417
  \end{itemize}
  \end{minipage}
  \begin{minipage}{0.34\textwidth}
     \centering
     \includegraphics[width=\textwidth]{Pictures/wenqing.pdf}
  \end{minipage}

  \item[$\bullet$] In both cases, please show that this established method worked well in the current data set: the Cu$+$Au $R_\gamma$ and direct photon spectrum are consistent with the Au$+$Au 40--50\% result from PPG243, which has comparable charged-particle multiplicity

  \item[$\bullet$] If you modified an established method, please highlight the modification you made: none


\end{itemize}
\end{frame}


\begin{frame}{Outline of the analysis method}
   \small
   The direct photon double ratio is measured as
   \begin{equation*}
      R_\gamma = \frac{\gamma^{\rm inclusive}}{\gamma^{\rm decay}}
      \approx
      \frac{
      \left\langle\varepsilon_\gamma f\right\rangle
      \left(N_{\rm incl}/N_{\rm tag}\right)_{\rm data}
      }{
      \left(\gamma^{h}/\gamma^{\pi^0}\right)_{\rm MC}
      } .
   \end{equation*}
   \begin{itemize}
      \item $N_{\rm incl}$ is the inclusive conversion-photon yield from the $M_{e^+e^-}$ distribution.
      \item $N_{\rm tag}$ is obtained by combining a conversion photon with an EMCal photon and extracting the $\pi^0$ peak in $M_{e^+e^-\gamma}$.
      \item $\left\langle\varepsilon_\gamma f\right\rangle$ is the probability to tag the partner photon from a $\pi^0$ decay, evaluated with simulation and occupancy correction.
      \item $\gamma^h/\gamma^{\pi^0}$ is the decay-photon cocktail correction from HELIOS.
   \end{itemize}
\end{frame}

\begin{frame}{Extraction of $N_{\rm incl}$ and $N_{\rm tag}$}
   \small
   \begin{columns}[T,onlytextwidth]
      \column{0.50\textwidth}
      \textbf{Inclusive conversion yield}
      \begin{itemize}
         \item Reconstruct external conversions, $\gamma\rightarrow e^+e^-$.
         \item Extract $N_{\rm incl}$ from the $M_{e^+e^-}$ distribution after background subtraction.
         \item The conversion mass shape is shifted because the tracks are reconstructed assuming the event vertex, while the conversion occurs in detector material at finite radius.
      \end{itemize}
      \column{0.50\textwidth}
      \textbf{Tagged $\pi^0$ yield}
      \begin{itemize}
         \item Combine the conversion photon with EMCal photon candidates.
         \item Extract the $\pi^0$ signal from the $M_{e^+e^-\gamma}$ distribution.
         \item Mixed-event background is normalized outside the $\pi^0$ peak and subtracted.
         \item Residual background is described by a polynomial under the $\pi^0$ peak.
      \end{itemize}
   \end{columns}
   \vspace{0.4em}
   The nominal $\pi^0$ extraction uses a Gaussian signal plus a second-order polynomial residual background.  Variations of the integration window, fit range, and polynomial order are included in the systematic uncertainty.
\end{frame}

   \begin{frame}{Systematic uncertainties}
   \small
   Systematic uncertainties are evaluated by repeating the full extraction for physically motivated cut and fit variations.  Variations are grouped by source:
   \begin{itemize}
      \item conversion-pair selection;
      \item electron identification;
      \item electron-track momentum threshold;
      \item EMCal photon $E_{\rm core}$ threshold;
      \item EMCal shower-shape requirement;
      \item photon energy scale;
      \item $\pi^0$ peak extraction.
   \end{itemize}
   \vspace{0.4em}
   The systematic uncertainty is evaluated group-by-group from the spread of the variation relative to the nominal result.  The final band uses asymmetric up/down uncertainties with a baseline lower bound to avoid underestimates from statistical fluctuations in individual variation points.
\end{frame}

% Efficiency and $R_\gamma$ for all cuts plot to show evidence of underlying physics in the analysis.
\begin{frame}{Efficiency and $R_\gamma$ for all cuts}
   \centering
   \begin{columns}

      \column{0.50\textwidth}
      \includegraphics[width=0.99\textwidth, trim=0 0 0 0.1cm, clip]{Pictures/ef_configs.png}
      \column{0.50\textwidth}
      \includegraphics[width=0.99\textwidth, trim=0 0 0 0.1cm, clip]{Pictures/rgamma_all_configs.png}
   \end{columns}
   \begin{itemize}
      \item There are clear trends in the efficiency as a function of the applied cuts.
      \item $R_\gamma$ shows no dependence on the applied cuts and counting windows.
      \item $R_\gamma$ agrees with 40-50\% Au+Au results at similar multiplicity to minimum bias Cu+Au.
   \end{itemize}
\end{frame}



\begin{frame}{Analysis Issues}
   \textbf{What were the main issues of the analysis?} \break
   \begin{itemize}
      \item[$\bullet$] The earlier $R_\gamma$ result was too flat as a function of $p_T$
      \item[$\bullet$] At high $p_T$, the simulated $\pi^0$ peak was too narrow because of the simulation vertex treatment, which caused a loss of tagging efficiency
      \item[$\bullet$] At low $p_T$, the ${n_{cluster} == 0}$ rejection removed inclusive conversion photons without a corresponding tagged photon contribution
      \item[$\bullet$] This made $\langle\varepsilon_\gamma f\rangle = N_{\rm tag}/N_{\rm incl}$ too large at low $p_T$
   \end{itemize}

   \textbf{What concerns were raised in PWG meetings?}
   \begin{itemize}
      %\item[$\bullet$] Stability of the tagging efficiency correction
      \item[$\bullet$] Shape of $R_\gamma$ compared to the Au$+$Au reference
      \item[$\bullet$] Treatment of systematic uncertainties
   \end{itemize}

   \textbf{What was the resolution for these issues?}
   \begin{itemize}
      \item[$\bullet$] Corrected the simulation vertex treatment, improving the high-$p_T$ $\pi^0$ peak and tagging efficiency
      \item[$\bullet$] Removed the ${n_{cluster} == 0}$ rejection to avoid losing inclusive photons at low $p_T$
      \item[$\bullet$] Re-evaluated $N_{\rm incl}$, $N_{\rm tag}$, $\langle\varepsilon_\gamma f\rangle$, $R_\gamma$, and the direct photon spectrum
      \item[$\bullet$] Estimated systematics by varying conversion cuts, electron ID, photon cuts, energy scale, and $\pi^0$ extraction settings
   \end{itemize}
\end{frame}

\begin{frame}{Relations to Previous Analyses}
   \textbf{Previously published PHENIX results for the same observable:}
   \begin{itemize}
      \item[$\bullet$] No published PHENIX Cu$+$Au direct photon result with the external conversion method
   \end{itemize}

   \textbf{Previously published PHENIX results with a similar observable and method:}
   \begin{itemize}
      \item[$\bullet$] PHENIX Au$+$Au direct photon measurement using the external conversion method
   \end{itemize}

   \begin{center}
      \includegraphics[width=0.45\textwidth]{Pictures/wenqing.pdf}
   \end{center}

   \textbf{Previously published non-PHENIX results for similar observables:}
   \begin{itemize}
      \item[$\bullet$] Yes, STAR and ALICE have published related direct photon measurements
   \end{itemize}
\end{frame}



\begin{frame}{Plots Requested for PHENIX Preliminary: $R_\gamma$}
   \centering
   \includegraphics[width=0.78\textwidth]{Pictures/rgamma_stat_syst0.png}
   \vspace{0.2em}

   \small Cu$+$Au minimum-bias $R_\gamma$ compared with the PHENIX Au$+$Au 40--50\% reference from PPG243.
\end{frame}

\begin{frame}{Plots Requested for PHENIX Preliminary: $R_\gamma$ multiplicity comparison}
   \centering
   \includegraphics[width=0.78\textwidth]{Pictures/rgamma_stat_syst1.png}
   \vspace{0.2em}

   \small Cu$+$Au minimum-bias $R_\gamma$ compared with PHENIX Au$+$Au 50--60\%, 40--50\%, and 30--40\% reference measurements from PPG243.
\end{frame}

\begin{frame}{Plots Requested for PHENIX Preliminary: direct photon spectrum}
   \centering
   \includegraphics[width=0.78\textwidth]{Pictures/spectrum_stat_syst.png}
   \vspace{0.2em}

   \small Direct photon invariant yield in Cu$+$Au minimum-bias collisions compared with the PHENIX Au$+$Au 40--50\% reference.
\end{frame}

\begin{frame}{Plots Requested for PHENIX Preliminary: integrated yield}
   \centering
   \includegraphics[width=0.7\textwidth]{Pictures/dndy_vs_dNch.png}
   \vspace{0.2em}

   \small Integrated direct photon yield $dN^{\rm dir}_{\gamma}/dy$ for $1<p_T<5$ GeV/$c$ versus charged-particle multiplicity at midrapidity.
\end{frame}

\begin{frame}{Physics message}
   \normalsize
   \begin{itemize}
      \item The measured Cu$+$Au minimum-bias $R_\gamma$ is consistent with the PHENIX Au$+$Au reference measurements at similar charged-particle multiplicity.
      \item The extracted direct photon invariant yield agrees with the Au$+$Au 40--50\% result within uncertainties.
      \item The Cu$+$Au integrated direct photon yield follows the global scaling of $dN^{\rm dir}_{\gamma}/dy$ with $dN_{\rm ch}/d\eta$ observed across collision systems and beam energies.
      \item The result extends the PHENIX external-conversion direct photon program to the asymmetric Cu$+$Au system.
   \end{itemize}
\end{frame}

\begin{frame}{The Path to the Final Results}
   \textbf{What improvements are expected for the final results?}\break
   Using whole available data and centrality bins.\break
   \break
   \textbf{Is all data available used?}\break
   No, about 30\% of the Run-12 Cu$+$Au minimum-bias sample is used.\break
   \break
   \textbf{How much statistical/systematical gain is expected before final publication?}\break
   About $\sqrt{3}$ statistical gain is expected. \break
   \break
   \textbf{Provide a detailed timeline to publication:}\break
   Prepare article draft by fall 2026.\break
   \break
   \textbf{Provide details on who will do the work and milestones towards final results:}\break
   Tongzhou Guo, A. Drees, and Yu. Mitrankov.
\end{frame}

\end{document}
